% Ch.10 Conclusion (~3-4 pp). Draft on Day 12.
\chapter{Conclusion}
\label{ch:conclusion}

This thesis asked one question: does the macro-anchored cycle factor of
\citet{cieslak2015}---the predictor of US bond returns with the strongest claim
to an economic mechanism---apply to international government bonds? To answer it,
the analysis built a three-factor hierarchy on a monthly panel of eleven G10
markets spanning 1989/1994--2026: the \emph{local} cycle factor $\CF_{i,t}$,
constructed for each market exactly as in the original study; the GDP-weighted
\emph{global} cycle factor $\GCF_{t}$; and a novel \emph{FX-adjusted} global
cycle factor $\FXGCF_{t}$ for an unhedged US-dollar investor. After validating
the estimation engine against the published US and international results
(\Cref{ch:replication}), the three factors were tested in three phases, in
sample and under a fully-recursive out-of-sample protocol
(\Cref{ch:results,ch:robustness,ch:strategy}). The answer is affirmative, but
with an important refinement: the cycle mechanism is a property of the
international asset class, yet the form of it that an investor can actually use
is the \emph{global} factor, not the local one.

\section{Answers to the Research Questions}
\label{sec:concl-answers}

\paragraph{The main question.} The \citet{cieslak2015} cycle factor does apply
to international government bonds. The local cycle factor predicts one-year
excess returns with exactly the US sign structure---a negative loading on the
one-year cycle, a positive loading on the average cycle---in all eleven markets,
significantly so in ten, with an average in-sample $R^{2}$ of $25\%$ (reaching
$33\%$ for the United States, closely matching the original) and loadings that
are flat-to-rising across the maturity curve. These magnitudes sit far above the
$95$th percentile of the predictive $R^{2}$ that the expectations hypothesis can
generate in a sample of this length (\Cref{sec:meth-ehmc}). The trend--cycle
mechanism is therefore not a US idiosyncrasy: in every G10 market, the deviation
of yields from their slow-moving inflation anchor measures the compensation
demanded for holding duration.

\paragraph{Sub-question (1): local versus a GDP-weighted global factor.} The
predictable component of international bond returns is overwhelmingly global. The
GDP-weighted global cycle factor $\GCF_{t}$ is statistically significant in
\emph{every} market and, on its own, explains as much return variation as the
eleven separately estimated local factors (a cross-country average $R^{2}$ of
$25\%$ for each). A single international factor thus reproduces the
predictability of eleven national ones---the integration result of
\citet{dahlquist2013}, here re-established for the macro-anchored cycle family
rather than the forward-rate family on which it was first documented.

\paragraph{Sub-question (2): does the global factor subsume the local factor?}
Largely, but not completely. Once the global factor is included, the
orthogonalised local factor retains significant incremental content in only
three of the eleven markets---Italy, the Netherlands, and Belgium---and is
indistinguishable from zero in the other eight. The exceptions are informative
rather than awkward. The Italian case decomposes cleanly: its local
predictability is overwhelmingly a 2010--2012 sovereign-debt-crisis phenomenon,
with a horse-race $t$-statistic above ten and the single largest positive
real-time local result in the study ($R^{2}_{\mathrm{oos}}=+0.21$), while the
global factor enters with the wrong sign (\Cref{sec:rob-italy}). Integration is
the rule; it bends exactly when redenomination and sovereign-credit risk reprice
a peripheral market that a core-dominated global factor cannot capture.

\paragraph{The dollar investor and the FX adjustment.} Currency risk is the
binding constraint for an unhedged US investor. Dollar conversion erodes roughly
two-thirds of the global factor's in-sample predictive power, and the
FX-adjusted factor restores only part of it: reliably in direction and across
markets, but modestly in magnitude, because in this G10 sample $\FXGCF_{t}$ is
nearly collinear with the unadjusted global factor (correlation $0.99$, against
about $0.50$ for the forward-rate analogue of \citet{dahlquist2013}). That the
FX adjustment which transforms the global \emph{forward} factor adds little to
the global \emph{cycle} factor is itself a finding: the detrending that gives the
cycle factor its real-time stability in bonds strips out precisely the
rate-level information that an FX adjustment needs to forecast currencies.

\paragraph{The decisive test: real time.} The conclusion that matters for
practice is the out-of-sample one. Of the four candidate factors, only the
global cycle factor beats the recursive prevailing mean: its pooled
$R^{2}_{\mathrm{oos}}$ is $+0.05$, positive in ten of eleven markets, while the
local factor ($-0.08$) and the unadjusted dollar-investor factor ($-0.07$) lose
to the historical mean almost everywhere. This real-time edge survives in every
subsample except the 2008 panic, beats the global forward-rate factor that fits
\emph{better} in sample, and depends materially on measuring the inflation
anchor with \emph{core} rather than headline CPI. It is also economically
exploitable: a currency-hedged timing strategy built on the global factor raises
the annual Sharpe ratio of a G10 bond portfolio from $0.31$ to $0.38$, with a
certainty-equivalent gain of about half a percent a year, earned chiefly by
sidestepping drawdowns such as 2022 (\Cref{ch:strategy}). The gap between the
local factor's strong in-sample fit and its negative out-of-sample $R^{2}$ is
the clearest sign that the predictability worth taking seriously is the global
one.

\section{Contribution}
\label{sec:concl-contribution}

The thesis contributes a three-factor hierarchy and the evidence to evaluate it.
It is, to the best of my knowledge, the first study to carry the
\citet{cieslak2015} cycle-factor \emph{construction}---rather than raw inflation
trends added to yield-curve regressions, as in \citet{zhang2021}---across the
full G10 cross-section, including the euro-area periphery whose sovereign crisis
provides a natural experiment on the limits of integration. It transplants the
integration test of \citet{dahlquist2013} from the forward-rate factor to the
macro-anchored cycle, and finds the integration result \emph{stronger} for the
cycle factor in real time. And it introduces the FX-adjusted global cycle factor,
documenting the substantive contrast with its forward-rate predecessor. Threaded
through all three phases is a methodological point: because the cycle factors are
generated regressors, the entire construction---trend, cycles, predictive
coefficients, and global weights---was re-estimated fully recursively and scored
with the \citet{campbellthompson2008} out-of-sample $R^{2}$. Applying that
discipline to the whole generated-regressor chain, not only to the final
regression, is what reverses the in-sample verdict and isolates the one robust
object the thesis stands on.

The honest one-line summary is therefore not ``the cycle factor works
everywhere'' but something more specific and more useful: international bond risk
premia contain a common, macro-anchored, slowly-moving component that is large
enough and stable enough to estimate in real time---provided one aggregates
across markets, measures the inflation anchor cleanly with core CPI, and does not
ask it to forecast exchange rates.

\section{Limitations and Future Work}
\label{sec:concl-future}

The conclusions are bounded by the limitations set out in \Cref{sec:disc-limitations}
and should be read with them in view: HAC inference can over-reject for
persistent regressors under overlapping returns; the sample spans one secular
disinflation and effectively far fewer than eleven independent ``global''
observations, given the shared euro-area monetary policy; the maturity menu and
the top-down FX construction are dictated by the data; the expectations-hypothesis
Monte Carlo reproduces the shape but not the exact level of the published null
band; and the global factor's real-time edge is a property of expanding-window,
core-inflation estimation that does not survive a rolling window or the 2008
crisis.

Several of these limitations point directly to future work. The most natural
extension is to embed the global cycle in a no-arbitrage term-structure model
with a shared international trend, closing the gap between the reduced-form
evidence assembled here and the structural ``falling-stars'' framework of
\citet{bauer2020}. A richer maturity grid and a longer, more balanced
cross-section---representing the euro area by a single benchmark and adding
currency-diverse markets such as Australia, New Zealand, and Norway, as in
\citet{randl2025}---would sharpen both the integration test and the FX
adjustment, and would let alternative weighting schemes (equal,
market-capitalisation, principal-component) be compared against the GDP weights
used here. Finally, the real-time global cycle factor is a candidate input to
cross-sectionally optimised, rather than purely timing, portfolios: the
out-of-sample evidence does not support country selection on the local factors,
but a forecast with the demonstrated stability of the global factor is exactly
the expected-return input that mean--variance-efficient international bond
portfolios are most sensitive to. Pursuing these directions would test how far
the single robust object identified here---a common, macro-anchored, real-time
global bond cycle---can be pushed.
